\documentclass[10pt,a4paper]{amsart}
\usepackage[margin=19mm]{geometry}
\usepackage{amsmath,amssymb,mathtools}
\usepackage[scr=rsfs,cal=euler]{mathalfa}
\usepackage{xfrac}
\usepackage[unicode,hidelinks,bookmarks=false]{hyperref}
\newcommand{\ie}{i.e.\ }
\newcommand{\cf}{cf.\ }
\newcommand{\resp}{respectively\ }
\newcommand{\Subsection}{\subsection}
\providecommand{\nobreakdash}{\nobreak\hbox{-}}
\setlength{\emergencystretch}{2em}
\multlinegap0pt
\usepackage{algorithm}\usepackage{algpseudocode}
\newtheorem{theorem}{Theorem}
\def\Det{{\rm det}}
\def\Tr{{\rm Tr}}
\title{On efficient approximation of quadratic irrationals, signed Chebyshev polynomials and Householder's method}
\author{Peter H. van der Kamp, Anthony Overmars, Marcel Jackson, Andrew N.W. Hone}
\date{Source: arXiv:2511.21617v5 (CC BY 4.0)}
\begin{document}
\maketitle

\noindent\emph{Source and licence.} On efficient approximation of quadratic irrationals, signed Chebyshev polynomials and Householder's method. Version arXiv:2511.21617v5 (CC BY 4.0), distributed by arXiv under the Creative Commons Attribution 4.0 International licence, http://creativecommons.org/licenses/by/4.0/, as recorded in the arXiv licence metadata for this exact version and shown on the abstract page https://arxiv.org/abs/2511.21617v5. Full licence notice: \texttt{licenses/arxiv-2511.21617-CC-BY-4.0.txt}.\par\smallskip

\noindent\textit{Editorial scope.} Counted unit: the article's theorem thm:tcs (source0.tex lines 679--687). The article's complete original proof of it is delivered with it (source0.tex lines 688--710, one \texttt{proof} environment of 23 lines). No mathematics is written, repaired or re-derived: the statement and the proof are the article's own lines, in the article's own notation, and no retained line is substituted or re-flowed.  Retained uncounted as the source-local context the proof needs: lines 122--124; lines 126--128; lines 229--232; lines 314--318.  Nothing was omitted from any retained span, so the package carries no omission record.  No retained line contains a citation, so the wrapper carries no bibliography and no bibliography entry.  No retained line contains a figure environment, an image-inclusion command or drawing code, so no image is delivered and the package declares no asset dependency.  Editorial formatting only, in addition: an AMS \texttt{amsart} preamble that restates the article's own declarations and macros used by the retained lines, namely the environments \texttt{theorem} on separate counters, together with the standard packages those lines need.  The retained environments print in this document as Theorem 1 in order of appearance, whereas the article prints the same items with the numbering of its own full document.  The complete native source of the article as distributed in the arXiv e-print for this exact version is bundled unchanged as \texttt{sources/189887/source0.tex}.
\begin{equation} \label{eq:sn}
\sqrt{N}=[c_0,\overline{c_1,c_2,\ldots,c_{2},c_{1},2c_0}]
\end{equation}

\begin{equation} \label{eq:Pll}
p_{l-1}^2-Nq_{l-1}^2=(-1)^l.
\end{equation}

\begin{align}
T^l_{k+2}&=2xT^l_{k+1}-(-1)^lT^l_k,\qquad T^l_0=1,\ T^l_1=x,\label{eq:Tl}\\
U^l_{k+2}&=2xU^l_{k+1}-(-1)^lU^l_k,\qquad U^l_0=1,\ U^l_1=2x.\label{eq:Ul}
\end{align}

\begin{align}
\Psi_{n+2l}&=\Psi_n\left(\Psi_n^{-1}\Psi_{n+l}\right)^2\notag\\
&=\Psi_n\left(\Tr(\Psi_n^{-1}\Psi_{n+l})\Psi_n^{-1}\Psi_{n+l}-\Det(\Psi_n^{-1}\Psi_{n+l})\right)\notag\\
&=\Tr(\Psi_n^{-1}\Psi_{n+l})\Psi_{n+l}-(-1)^l\Psi_{n},\label{eq:psir}
\end{align}

\begin{theorem} \label{thm:tcs}
Let $p_i,q_i$, $i=0,1,\ldots$, be the numerators and denominators of the convergents of a continued fraction of the form \eqref{eq:sn}, so that equation \eqref{eq:Pll} holds. We then have
\begin{equation} \label{eq:pqC}
\begin{split}
p_{kl-1}&=T^l_k(p_{l-1}),\\
q_{kl-1}&=q_{l-1}U^l_{k-1}(p_{l-1}).
\end{split}
\end{equation}
\end{theorem}

\begin{proof}
The palindromic nature of \eqref{eq:sn} tells us that the matrix $\Psi_0^{-1}\Psi_{l-1}$ is symmetric, which yields $q_{l-2}=p_{l-1}-c_0q_{l-1}$. Together with the fact that
$c_l=2c_0$, this implies that
\begin{align*}
t_1&=\Tr\left(\Psi_0^{-1}\Psi_{l}\right)\\
&=\Tr\left(
\begin{pmatrix} 0 & 1 \\ 1 & -c_0 \end{pmatrix}
\begin{pmatrix} p_{l-1} & p_{l-2} \\ q_{l-1} & q_{l-2} \end{pmatrix}
\begin{pmatrix} 2c_0 & 1 \\ 1 & 0 \end{pmatrix}\right)\\
&=c_0q_{l-1}+p_{l-1}+q_{l-2}\\
&=2p_{l-1}.
\end{align*}
The ${}_{1,1}$-component of formula \eqref{eq:psir} tells us that, with $p_{l-1}=x$, for all $n$,
\[
p_{n+2l}=2xp_{n+l}-(-1)^lp_{n}.
\]
Since $p_{-1}=1$, according to \eqref{eq:Tl} we have $p_{kl-1}=T^l_k(x)$.
Similarly, the ${}_{2,1}$-component of formula \eqref{eq:psir} yields
\[
q_{n+2l}=2xq_{n+l}-(-1)^lq_{n},
\]
and, as $q_{-1}=0$, this implies $q_{2l-1}=2xq_{l-1}$. Thus, we see that up to a constant factor $q_{l-1}$, the decimation $q_{kl-1}$ is described by a signed Chebyshev polynomial of the second kind, $U^l_{k-1}(x)$, cf. \eqref{eq:Ul}.
\end{proof}

\end{document}
